%!TEX root=../main.tex
\section{Related Work}
\label{sec:related_works}


\PHB{RL Post-Training Systems.}
%说一下现在的一些post-RL training systems 
Many systems address the systems challenges of RL post-training. Early frameworks~\cite{Harper_NeMo_a_toolkit,deepspeedchat,hu2024openrlhf} adopt static, stage-specific GPU partitions, which lowers overall utilization when the workload balance shifts across stages. veRL~\cite{verl} instead uses a hybrid controller that co-locates multiple stages on the same GPUs to improve hardware efficiency. Subsequent systems~\cite{liu2025specrlacceleratingonpolicyreinforcement,chen2025respecoptimizingspeculativedecoding,cheng2025fastllmposttrainingdecoupled,shao2025beatlongtaildistributionaware,RhymeRL,zhong2024rlhfuseefficientrlhftraining,areal,asyncflow,Laminar,StreamRL} explore a range of acceleration techniques, including speculative decoding, fusion of pipeline stages to reduce orchestration overhead, asynchronous or decoupled training to hide latency, and various forms of resource disaggregation. \SystemName{} builds on the disaggregated and asynchronous paradigm and assigns workloads based on hardware affinity.



\PHM{Resource Disaggregation.} The concept of resource disaggregation, first proposed in early architectural work on memory~\cite{disaggregated-mem-isca09}, has become a cornerstone of modern system design. Recent systems like LegoOS~\cite{shan2018legoos} and Mira~\cite{guo2023mira} decouple hardware into separate resource pools that can be managed independently to improve utilization. This pattern is prevalent in LLM serving, where systems commonly disaggregate the prefill and decoding phases~\cite{distserve, patel2023splitwise, hu2024inference, strati2024dejavu, qin2024mooncake}, with some even disaggregating transformer sub-modules like attention and FFNs~\cite{megascale-infer,step3}. In the context of RL post-training, several works apply a similar principle, separating the training and rollout stages across different GPU types~\cite{StreamRL,seamlessflow}. Compared to these approaches, \SystemName{} provides a more general and fine-grained disaggregation model, tailored specifically for the entire lifecycle of multi-task agentic RL. 

